% !TEX root = ../thesis.tex
%
% cover page Medieninformatik
% @author Thomas Lehmann
% bases on work by unknown
%

\thispagestyle{empty}
% In dieser Umgebung wird die Titelseite separat vom restlichen Text gesetzt
\begin{titlepage}
  % andere Seitenränder als im Rest der Arbeit
  \newgeometry{lmargin=2cm,tmargin=7mm,rmargin=5mm,bmargin=1cm}
  % die "Hausfarbe" der HAW; diese und die folgenden Einstellungen sind lokal
  % und gelten nur innerhalb der Umgebung "titlepage"
  %\color{haw}
  % Blocksatz für die Titelseite deaktivieren
  \raggedright
  % Logo rechtsbündig setzen
  \hfill\includegraphics[width=7cm]{../style/HAW_Marke_RGB_300dpi}\\

  % vertikaler Abstand
  \vspace{5cm}

  % Wahl der "Hausschrift" Open Sans der HAW, die als Schrift auf Ihrem
  % Rechner installiert sein muss
  %\setmainfont{Open Sans}
  % etwas kleiner als üblich
  {\fontfamily{phv}\selectfont
  \fontsize{10pt}{10pt}
  % fett und in Majuskeln
  \textcolor{haw}{\textbf{BACHELORARBEIT}}
  }
  % vertikaler Abstand
  \vspace{8mm}

  % der Titel der Arbeit als "Seite in der Seite"; natürlich müssen Sie hier
  % Ihren Titel eintragen
  \begin{minipage}{0.8\linewidth}
    % Wahle der zweiten "Hausschrift" der HAW, die ebenfalls auf Ihrem Rechner
    % bereits vorhanden sein muss
    %\setmainfont{Martel Heavy}
    % ziemlich große Schrift
    {
    \fontfamily{cmr}
    \ICDTitleFontSize
    \selectfont
    % [1mm] steht jeweils für einen etwas größeren Durchschuss
    \textcolor{haw}{\textbf{\IthesisTitle}}\\
    % am Ende noch ein waagerechter Strich, das CD will es so...
    \textcolor{haw}{\,\rule{11mm}{1.2mm}}
    }
  \end{minipage}

  % vertikaler Abstand, überraschenderweise
  \vspace{1cm}
  
  {\fontfamily{phv}\selectfont
  \fontsize{10pt}{10pt}
  
  % hier korrektes Datum und Ihren Namen eingeben
  \textcolor{haw}{vorgelegt am \ISubDate \\
  \IthesisAuthor}

  % letzter vertikaler Abstand für heute
  \vspace{5cm}

  % noch eine "Seite in der Seite", etwas nach rechts geschoben
  \hspace*{37mm}
  \begin{minipage}{0.5\linewidth}
    % Namen und Titel der beiden Prüfer eintragen
    \textcolor{haw}{
    \begin{tabular}{@{}ll}
      Erstprüferin: & \IfirstSv\\[-.3mm]
      Zweitprüfer: & \IsecondSv \\
    \end{tabular}\\
    }
    % noch ein horizontaler Strich
    \textcolor{haw}{\,\rule{9mm}{1mm}\\[1.5mm]}

    \textcolor{haw}{
    \textbf{HOCHSCHULE FÜR ANGEWANDTE}\\
    \textbf{WISSENSCHAFTEN HAMBURG}\\
    Fakultät Informatik und Digitale Gesellschaft\\
    Finkenau 35\\
    22081 Hamburg}
  \end{minipage}
  }
  
  % setzt die Geometrie wieder auf die Standardwerte zurück
  \restoregeometry
\end{titlepage}
\cleardoublepage